% ============================================================
%  Title: Rombotics Q8: An Open-Source Sub-£250 Quadruped Robot
%         with Embedded Per-Axis PID Balance Control,
%         Loop-Timing Characterisation, and Independent Sensor Validation
%
%  Author : Romil Chetankumar Modi
%           MEng Mechatronics and Robotics Engineering (Distinction)
%           Anglia Ruskin University, Cambridge, UK
%           rombotics.com | github.com/romil0001
%
%  Supervisors: Dr Shabnam Sadeghi Esfahlani,
%               Dr Hossein Bisheh
%
%  Compile : pdflatex -> bibtex -> pdflatex x2
%
%  v7 Changelog (reviewer-response pass — Round 1):
%  -------------------------------------------------------
%  MAJOR FIXES
%  M1 — Abstract MAE now reports both pitch (0.18°) and roll (0.21°);
%       selective single-value reporting corrected throughout.
%  M2 — "Independent validation" claim scoped to pitch/roll only in
%       abstract, Contributions §I.A, §V-B, and conclusion; yaw
%       inclinometer limitation stated explicitly in each location.
%  M3 — Yaw 20% improvement baseline now given with n=10, mean±CI,
%       and a two-sample t-test result (t(38)=6.36, p<0.001) in §V-D.
%  M4 — t-test statistics added: yaw diagonal vs. uniform baseline
%       in §V-D; PID-on vs. PID-off gait comparison in §VI (t(6)=9.0,
%       p<0.001); two-sided α=0.05 as declared in §V-A.
%  M5 — Platform comparison table: "SpotMini EDU" (wrong citation)
%       corrected to "Stanford Doggo" with proper kau2019stanford ref;
%       body-text description in §II-D corrected accordingly.
%  M6 — Long-duration drift: abstract and §V-G now state "n=5;
%       preliminary characterisation" rather than presenting as a
%       fully-powered validated result.
%  M7 — Locomotion framing reconciled: abstract aligned with §VI
%       body text ("platform capability baseline" language used
%       consistently; "demonstrates … on level ground" qualified).
%
%  MINOR FIXES
%  m1 — Truncated Fig. 14 (MATLAB) caption trailing semicolon removed;
%       sentence completed.
%  m2 — MPU6050_tockn library commit hash and α value documented in
%       §IV-B; readers directed to pinned release for reproducibility.
%  m3 — Payload capacity methodology (static dead-weight test) added
%       to Table I footnote and §III-A.2.
%  m4 — PWM-space tibia error bound shown monotonic over ±10° range
%       via signed-deviation argument in §IV-E.
%  m5 — Busy-wait design choice justified vs. hardware-timer ISR in §IV-F.
%  m6 — ZMP analysis scope (static stance only; quasi-static criterion)
%       stated more explicitly in §III-C.
%  m7 — Ziegler–Nichols relay tuning parameters (Ku, Tu, ZN-initial vs.
%       final gains) added as Table V and explained in §IV-C.1.
%  m8 — Thermal characterisation note added to §V-J Discussion.
%
%  Additional clarifications (reviewer questions Q1–Q3):
%  Q1 — PWM-to-degree conversion factor documented in §III-B.4.
%  Q2 — Relationship between Table VIII (5° step) and Table IX
%       (10° step) clarified in §V-F.
%  Q3 — FreeRTOS background-task preemption note added in §IV-F.
% ============================================================
\documentclass[journal]{IEEEtran}
\usepackage[T1]{fontenc}
\usepackage[utf8]{inputenc}
\usepackage{amsmath,amssymb,amsfonts}
\usepackage{algorithm}
\usepackage{algorithmic}
\usepackage{graphicx}
\usepackage{booktabs}
\usepackage{array}
\usepackage{multirow}
\usepackage{cite}
\usepackage{url}
\usepackage{hyperref}
\usepackage{xcolor}
\usepackage{siunitx}
\usepackage{microtype}
\usepackage{balance}
\usepackage{subcaption}
\usepackage{dblfloatfix}   % Fixes double-column float placement
\usepackage{flafter}       % Floats never appear before their call-out in text
% Float placement tuning: prevent back-to-back table stacking
\setcounter{topnumber}{2}
\setcounter{bottomnumber}{2}
\setcounter{totalnumber}{4}
\renewcommand{\topfraction}{0.85}
\renewcommand{\bottomfraction}{0.85}
\renewcommand{\textfraction}{0.10}
\renewcommand{\floatpagefraction}{0.75}
\hyphenation{op-tical net-works semi-conduc-tor quad-ru-ped}
\hypersetup{colorlinks=true,linkcolor=blue,citecolor=blue,urlcolor=blue}
\newcommand{\vect}[1]{\boldsymbol{#1}}
\newcommand{\mat}[1]{\mathbf{#1}}
% Degree symbol — robust across all math/text contexts
\newcommand{\dg}{^{\circ}}
% ============================================================
\begin{document}
\title{Sim-to-real balance control of an 8-DoF quadruped using ROS 2 and embedded complementary filtering}
\author{Romil~Chetankumar~Modi%
\thanks{Supported by the Anglia Ruskin University (ARU) School of Engineering and the Built Environment. \emph{(Corresponding author: Romil C.~Modi.)}}%
\thanks{R.\,C.\,Modi is with ARU, Chelmsford, CM1\,1SQ, UK (e-mail: \href{mailto:romil2002modi@gmail.com}{romil2002modi@gmail.com}). This work stems from a May 2025 MEng Dissertation supervised by Dr.~S.\,S.~Esfahlani and Dr.~H.~Bisheh. Portions of Sections~III--IV were summarised in a departmental seminar at the ARU School of Engineering (July 2025); that event had no published proceedings and no submitted manuscript. Open data and firmware: \url{https://github.com/romil0001/Sim-to-Real_balance_control_of_an_8-DoF_quadruped_using_ROS-2_and_embedded_complementary_filtering}.}%
}
\maketitle
% ── Abstract ─────────────────────────────────────────────────
% [M1] Both pitch (0.18°) and roll (0.21°) MAE now reported.
% [M2] "Independent validation" scoped to pitch/roll; yaw limitation stated.
% [M6] Long-duration drift qualified as n=5 preliminary characterisation.
% [M7] Locomotion framing aligned with Section VI body text.
\begin{abstract}
Low-cost legged robot research is constrained by a reproducibility
gap: open-source platforms at or below \pounds300 either lack
onboard closed-loop balance, omit independent sensor validation,
or both.
This paper addresses that gap with the \textit{Rombotics Q8}, a
self-balancing quadruped actuated by eight high-torque servo motors
(two revolute joints per leg: femur and tibia), driven by an ESP32
dual-core microcontroller and a PCA9685 16-channel PWM driver.
Orientation is estimated at \SI{100}{\hertz} using an MPU6050
six-axis IMU fused through a complementary filter ($\alpha=0.98$).
Independent cross-validation against a Wixey WR300 digital inclinometer
confirms static balance offsets of $0.28 \pm 0.12\dg$ (pitch) and
$-0.33 \pm 0.08\dg$ (roll), with mean absolute errors of $0.18\dg$
(pitch) and $0.21\dg$ (roll) over 20 static trials; yaw was not
cross-validated as the inclinometer measures tilt angle only — an
external validation step absent from all reviewed open-source platforms
at this cost tier.
A per-axis PID controller with first-order output smoothing
compensates pitch, roll, and yaw disturbances: pitch and roll
recover from $\pm10\dg$ step inputs in $0.41 \pm 0.05\,\text{s}$,
yaw in $0.61 \pm 0.07\,\text{s}$.
A PWM-space piecewise-linear tibia-coupling law
couples each tibia servo inversely to its femur partner, ensuring
approximate four-point ground contact during balance corrections
without an additional joint-angle sensor; the approximation error is
analytically bounded to $<\SI{7}{\milli\metre}$ over the
$\pm10\dg$ balance operating range.
The \SI{10}{\milli\second} control loop executes with
$0.23\,\text{ms}$ jitter (max observed \SI{10.47}{\milli\second})
measured over \num{10000} consecutive iterations on physical hardware.
Long-duration testing ($n=5$ runs; preliminary characterisation)
indicates near-linear gyroscope drift of
$\SI{0.003}{\degree\per\minute}$ (pitch/roll) and
$\SI{0.0033}{\degree\per\minute}$ (yaw) over \SI{60}{\minute}.
Preliminary open-loop trot trials on a single flat surface establish
a directed-locomotion capability baseline at \SI{0.18}{\metre\per\second},
providing a foundation for future closed-loop gait development.
The complete platform costs under \pounds250 in components and
is fully open-source, providing a reproducible low-cost baseline
for legged-robot research.
\end{abstract}
\begin{IEEEkeywords}
Quadruped robot, self-balancing, PID control, complementary filter,
ESP32, MPU6050, PCA9685, inverse kinematics, servo mirroring,
legged locomotion, modular robotics, open-source robotics,
low-cost platforms.
\end{IEEEkeywords}
% ============================================================
\section{Introduction}
\label{sec:intro}
\IEEEPARstart{L}{egged} locomotion confers superior adaptability
over wheeled and tracked platforms in unstructured environments
such as disaster zones, industrial inspection sites, and
agricultural fields \cite{biswal2021development,pratt2000legged}.
Quadruped robots are particularly attractive because four contact
points provide an inherent static support polygon, and their
bio-inspired morphology enables traversal of complex terrain
inaccessible to wheeled vehicles \cite{bledt2018cheetah,hutter2016anymal}.
Yet the most capable contemporary platforms ---
Boston Dynamics Spot (\$74\,500), Unitree A1 (\$15\,000),
Unitree Go2 (\$2\,700) \cite{unitree2023go2} ---
remain financially inaccessible to most small research teams
and educational institutions, restricting reproducible
benchmarking and curriculum adoption.
Open-source alternatives exist at sub-\$500 cost points, but
a critical reproducibility gap persists: published platforms at
this tier either omit onboard closed-loop balance entirely
(relying on off-board compute or simulation only), or lack
independent sensor validation of their claimed performance.
This paper addresses that gap directly.
A central challenge in low-cost quadruped design is achieving
robust real-time self-balancing under dynamic disturbances using
computationally affordable embedded hardware without the benefit
of force-torque sensors, high-resolution encoders, or off-board
compute.
Classical stability theory anchors this problem: the robot must
continuously maintain its Centre of Gravity (CoG) within the
support polygon defined by its four feet
\cite{vukobratovic2004zmp,jia2018stability}, adjusting leg
positions in real time as the body tilts.
This paper describes the \textit{Rombotics Q8} --- a fully
3D-printed, 8-servo quadruped that achieves robust closed-loop
self-balancing entirely on a single ESP32 at a total build cost
under \pounds250 ($\approx$\$315).
To the authors' knowledge, no published platform at this cost
point simultaneously provides: (i) an active onboard closed-loop
three-axis balance controller, (ii) external independent sensor
validation of pitch and roll orientation, and (iii) fully
open-source firmware with a documented loop-timing budget.
The Q8 is designed to fill that gap and to serve as a reproducible
benchmark substrate for future embedded reinforcement learning and
sim-to-real transfer research.
\noindent\textbf{Scope.}
This paper is scoped as a \emph{platform and balance-control
contribution}.
The primary validated claims are static and dynamic balance
performance on a level surface, control-loop timing, and
independent pitch/roll sensor cross-validation.
Directed locomotion (Section~\ref{sec:gait}) is presented as a
supplementary platform-capability baseline only; the open-loop
gait is not the primary validation target, and the limitations
of that result are stated explicitly.
Three aspects of the design are of particular technical
interest to the low-cost robotics community.
First, we develop and characterise a \emph{per-axis} PID
tuning strategy for embedded three-axis balance: pitch and roll
share gains reflecting their accelerometer-aided sensor fusion,
while yaw gains are independently tuned to account for its
higher moment of inertia and gyroscope-only estimation, a
nuance absent from most comparable open-source platforms.
Second, we present and characterise a \emph{PWM-space tibia-coupling
approximation} that avoids real-time trigonometry entirely, enabling
2-DoF leg compensation within the \SI{10}{\milli\second}
loop budget of a microcontroller, with foot-placement error
analytically bounded to $<\SI{7}{\milli\metre}$ over the
balance operating range.
Third, we provide a rigorous \emph{loop jitter characterisation}
over \num{10000} iterations and an \emph{independent inclinometer
cross-validation} of IMU pitch/roll accuracy --- experimental
rigour absent from all reviewed platforms at this cost tier
(Table~\ref{tab:compare}).
\begin{figure}[!t]
  \centering
  \includegraphics[width=\columnwidth]{fig11.jpg}
    \caption{Rombotics Q8 prototype in four operating postures.
    \textit{Top-left}: neutral four-point stance.
    \textit{Top-right}: front view showing femur/tibia articulation.
    \textit{Bottom row}: active balance correction sequences during
    external disturbance tests. White PLA structural links and orange
    TPU compliant joints are visible; the ESP32/PCA9685 electronics
    stack is mounted on the central chassis. Total build cost:
    under \pounds250.}
  \label{fig:robot_photo}
\end{figure}
\begin{figure*}[!t]
  \centering
  \includegraphics[width=\textwidth]{quadruped_robot_timeline.pdf}
  \caption{Historical milestones in quadruped robot development (1969--2020+).
    Key platforms include the GE Walking Truck (1969), the Adaptive Suspension
    Vehicle at Ohio State University (1980s), Boston Dynamics BigDog (2005),
    ETH Z\"{u}rich ANYmal (2015), Boston Dynamics LS3 and Spot (2015--2016),
    MIT Cheetah series (2018), Unitree Laikago (2017), ANYmal\,C (2019),
    and widespread commercial proliferation post-2020
    \cite{biswal2021development,chai2022survey,hutter2016anymal,bledt2018cheetah}.}
  \label{fig:timeline}
\end{figure*}
\subsection{Contributions}
\label{subsec:contributions}
\begin{enumerate}
% [M2] Scoped to "pitch and roll" for the cross-validation claim.
\item A modular, fully 3D-printed 8-DoF quadruped platform
      (Fig.~\ref{fig:robot_photo}) validated over 60-minute
      continuous balance trials with independent external
      pitch/roll sensor cross-validation, providing a
      reproducible low-cost baseline for legged-robot research.
% [M3][M4] Baseline now given with n, mean±CI, and t-test result.
\item A per-axis PID parameter set with physical rationale:
      pitch/roll ($K_p=1.5, K_i=0.1, K_d=0.2$) and yaw
      ($K_p=1.2, K_i=0.05, K_d=0.35$) with a first-order
      smoother ($\alpha_s=0.5$) eliminating servo chatter,
      and a derivation of the correct diagonal-differential
      yaw actuation geometry for four-legged platforms,
      yielding a 20\% improvement in yaw recovery time
      relative to a uniform-sign baseline
      ($n=10$ reference trials, $0.75\pm0.08\,\text{s}$;
      two-sample $t$-test: $t(38)=6.36$, $p<0.001$).
\item A static stability analysis applying the ZMP criterion
      to the Q8 geometry, establishing a $\pm49$--$53\dg$
      static tipping margin and motivating the balance
      controller's corrective role within that envelope.
\item A PWM-space piecewise-linear tibia-coupling approximation
      that reduces per-leg compute from \SI{31}{\micro\second}
      to $<\SI{1}{\micro\second}$, with analytically bounded
      foot-placement error $<\SI{7}{\milli\metre}$ over the
      $\pm10\dg$ balance operating range.
\item A quantitative loop-timing characterisation:
      \SI{10}{\milli\second} target, $\sigma=0.23\,\text{ms}$,
      max observed \SI{10.47}{\milli\second} over \num{10000}
      iterations on physical hardware.
% [M2] Scoped: validation applies to pitch/roll; yaw deferred.
\item A complementary-filter drift characterisation:
      $\SI{0.003}{\dg\per\minute}$ (pitch/roll),
      $\SI{0.0033}{\dg\per\minute}$ (yaw); pitch/roll
      independently cross-validated against a Wixey WR300
      inclinometer (yaw heading validation deferred to
      future magnetometer integration; see Section~\ref{subsec:drift}).
\item Preliminary open-loop trot gait results at
      \SI{0.18}{\metre\per\second} on a single flat surface,
      establishing a directed-locomotion capability baseline
      with explicit quantification of PID-enabled stance
      stability, as a foundation for closed-loop gait.
\item Fully open-source firmware (MajorQuad\_v0\_3) at
      \url{https://github.com/romil0001}.
\end{enumerate}
\subsection{Paper Structure}
Section~\ref{sec:related} reviews related work.
Section~\ref{sec:hw} describes the hardware design including
a ZMP static stability analysis.
Section~\ref{sec:sw} presents the software and control
architecture.
Section~\ref{sec:results} reports experimental results.
Section~\ref{sec:gait} describes preliminary gait trials.
Section~\ref{sec:conc} concludes.
% ============================================================
\section{Related Work}
\label{sec:related}
\subsection{Self-Stabilisation in Legged Robots}
Self-stabilisation requires continuous adjustment of the CoG
relative to the support polygon.
The Zero Moment Point (ZMP) criterion \cite{vukobratovic2004zmp}
is a classical sufficient condition for stable legged locomotion;
maintaining ZMP inside the support polygon yielded 90\% success
in post-disturbance recovery on the HyQ platform
\cite{chai2022survey}.
More general metrics such as the Gravito-Inertial Wrench Cone
(GIWC) \cite{caron2015stability} extend stability analysis to
uneven surfaces.
For low-cost platforms lacking foot-force sensors, body-angle
control via IMU feedback provides a pragmatic proxy:
Jia et al.\ \cite{jia2018stability} report a dynamic gait
stability criterion based solely on body angular velocity,
directly applicable to sensor-sparse designs such as the Q8.
High-performance platforms address stability through whole-body
model-predictive control (MPC): Di Carlo et al.\
\cite{dicarlo2018mpc} demonstrated high-bandwidth CoM regulation
on MIT Cheetah~3 using a convex MPC formulation solved at
\SI{500}{\hertz}. Gehring et al.\ \cite{gehring2016anymal}
achieved robust posture control on ANYmal through
hierarchical QP-based whole-body control.
While these methods achieve superior performance, their
computational demands (\si{\giga\flops}-class hardware,
Pinocchio/RBDL libraries) make them unsuitable for a
\SI{240}{\mega\hertz} microcontroller.
\begin{figure}[!t]
  \centering
  \includegraphics[width=1\linewidth]{fig16.png}
  \caption{Generalised multi-link humanoid kinematic model illustrating
    joint-angle parameterisation, body coordinate frames
    ($O_T$, $O_B$, $O_F$), and the Zero Moment Point (ZMP) reference
    \cite{biswal2021development}. The high dimensionality of whole-body
    control on such platforms ($>30$ DoF) motivates the simplified
    8-DoF design of the Q8 to reduce embedded compute demands.}
  \label{fig:kinematic_ref}
\end{figure}
\subsection{Learning-Based Locomotion}
Reinforcement learning (RL) has produced compelling
locomotion results on high-end platforms.
Kumar et al.\ \cite{kumar2021rma} demonstrated rapid motor
adaptation (RMA) on Unitree A1 using two-stage policy
distillation from simulation.
Lee et al.\ \cite{lee2020learning} achieved agile locomotion
on ANYmal across challenging terrains using a
teacher-student RL framework.
Rudin et al.\ \cite{rudin2022learning} demonstrated that
massively parallel GPU simulation in Isaac Gym can produce
deployable locomotion policies in under four minutes.
These results motivate future work on the Q8 (Section~\ref{sec:conc}),
as the platform's open-source structure and low rebuild cost
make it well-suited to sim-to-real transfer experiments.
Recent work has demonstrated deployment of distilled RL policies
onto microcontroller-class hardware via TensorFlow Lite Micro
\cite{hwangbo2019learning,rudin2022learning}, and the Q8's
documented timing budget (Section~\ref{subsec:timing}) provides
the headroom needed for lightweight policy inference at
\SI{100}{\hertz}.
\subsection{Sensor Fusion for IMU Orientation}
Accurate body-state estimation is foundational for balance
control.
Extended Kalman Filters (EKF) achieve $\sim\SI{0.5}{\dg}$
orientation accuracy but demand significant memory and
floating-point throughput
\cite{sabatini2011kalman,vitali2021robust}.
Madgwick filters offer an elegant gradient-descent alternative
with configurable convergence rate \cite{madgwick2010efficient},
achieving comparable accuracy at lower computational cost.
Complementary filters exploit the complementary frequency
responses of accelerometers and gyroscopes to achieve
$1$--$2\dg$ accuracy \cite{oz2022complementary} with
negligible overhead on constrained microcontrollers.
Since the ESP32 must simultaneously manage I$^2$C (two devices),
PWM dispatch (eight channels), and serial logging within a
\SI{10}{\milli\second} loop, the complementary filter embedded
in the MPU6050\_tockn library provides a well-matched
accuracy--speed trade-off.
% [M5] Corrected: platform is Stanford Doggo (kau2019stanford),
%      not SpotMini EDU. Description updated accordingly.
\subsection{Low-Cost Quadruped Platforms}
The \$350 Iron Dog Mini \cite{rahman2023irondogmini} demonstrated
Bézier-based gait generation with a 12-DoF leg design but relied
on PyBullet simulation without onboard closed-loop balance.
Stoch~2 \cite{tirumala2019gait} synthesised gaits via augmented
random search running on an off-board workstation with
hardware-in-the-loop.
The 12-DoF Spot Micro \cite{spotmicro2020} achieved trot gaits
at $\sim\SI{0.1}{\metre\per\second}$ using a Raspberry Pi,
but lacked real-time balance compensation.
The Stanford Doggo \cite{kau2019stanford} demonstrated
whole-body control on an open-source quasi-direct-drive
platform at \$3\,000, still an order of magnitude above
the Q8 cost.
A clear gap therefore exists in the published literature.
No open-source quadruped platform at or below \pounds300
(${\approx}$\,\$380) simultaneously provides active onboard
closed-loop balance, independent external sensor validation
(pitch/roll), and fully reproducible open-source firmware.
The Q8 targets all three.
\subsection{Embedded Control on Microcontrollers}
Chen et al.\ \cite{chen2021embedded} demonstrated that
ESP32-class hardware running FreeRTOS can sustain \SI{200}{\hertz}
control loops for aerial vehicles.
Liu et al.\ \cite{liu2024sensing} review sensing architectures
for indoor autonomous robots, noting that I$^2$C bus saturation
is the principal timing bottleneck in multi-sensor embedded
systems --- a finding directly relevant to our loop-jitter
analysis (Section~\ref{subsec:timing}).
Modular architecture principles from Jahn et al.\
\cite{jahn2019modular} informed the Q8's header-file-based
separation of servo limits, PID gains, and IK helpers.
% ============================================================
\section{Hardware Design}
\label{sec:hw}
\subsection{Mechanical Architecture}
The Q8 adopts a mammalian-inspired morphology with four identical
two-link legs (femur + tibia) attached symmetrically to a
rectangular base chassis (Fig.~\ref{fig:topdown}).
Each leg provides two revolute degrees of freedom: femur
flexion/extension and tibia flexion/extension, yielding 8~DoF
in total.
Hip abduction joints are deliberately omitted.
This 8-DoF limitation inherently restricts the robot's ability to
execute complex lateral manoeuvres or recover from severe lateral
dynamic disturbances compared to standard 12-DoF platforms.
However, removing the hip abduction DoF significantly simplifies
the control problem (2-DoF IK vs.\ 3-DoF IK), reduces the component
count by four servos (saving $\approx$\pounds48 and $\approx$250\,g),
and retains sufficient pitch and roll
compensation for foundational balance research at this cost tier.
Specifically, the 8-DoF design retains full sagittal-plane (pitch)
and frontal-plane (roll) balance authority, which are the dominant
disturbance axes for stationary or slow-moving legged robots on
level ground \cite{vukobratovic2004zmp}.
Lateral body translation and terrain-adaptive foot placement
--- the capabilities lost by omitting the hip DoF --- are
deferred to a future 12-DoF revision once the present
closed-loop balance baseline is established.
The deliberate trade-off is revisited in Section~\ref{sec:conc}
in the context of closed-loop gait integration.
Table~\ref{tab:dimensions} summarises key mechanical parameters.
% [m3] Payload methodology added as table footnote and in §III-A.2.
\begin{table}[!t]
\caption{Key Mechanical Parameters}
\label{tab:dimensions}
\centering
\begin{tabular}{lcc}
\toprule
\textbf{Parameter} & \textbf{Value} & \textbf{Unit} \\
\midrule
Chassis (L $\times$ W $\times$ H)   & $220\times180\times35$   & \si{\milli\metre} \\
Femur link length ($L_1$)           & 150                      & \si{\milli\metre} \\
Tibia link length ($L_2$)           & 150                      & \si{\milli\metre} \\
Standing height (neutral pose)      & 220                      & \si{\milli\metre} \\
Foot-to-foot span (lateral)         & 310                      & \si{\milli\metre} \\
Foot-to-foot span (longitudinal)    & 280                      & \si{\milli\metre} \\
Support polygon area (neutral)      & $310\times280$           & \si{\milli\metre\squared} \\
Total assembled mass                & 2.3                      & \si{\kilo\gram}   \\
Payload capacity (tested)$^{a}$     & 0.8                      & \si{\kilo\gram}   \\
Toe-out angle (from centreline)     & 45                       & \si{\degree}      \\
PLA infill                          & 40                       & \si{\percent}     \\
FDM layer height                    & 0.2                      & \si{\milli\metre} \\
\bottomrule
\multicolumn{3}{l}{\small Static dead-weight test: calibrated masses placed on} \\
\multicolumn{3}{l}{\small \phantom{$^{a}$\,}the chassis centroid; balance maintained for \SI{30}{\second}.} \\
\end{tabular}
\end{table}
\begin{figure}[!t]
  \centering
  \includegraphics[width=0.5\linewidth]{image.png}
  \caption{Top-down CAD view of the Q8 showing the rectangular
    chassis, leg attachment symmetry, and 45\textdegree{} toe-out
    stance. All four legs are geometrically identical;
    left/right-side role is assigned through servo polarity
    calibration in \texttt{ServoLimits.h}.}
  \label{fig:topdown}
\end{figure}
\subsubsection{Fabrication}
Structural components (base chassis, femur links, tibia links,
servo mounts) were designed in SolidWorks 2024 and fabricated
by FDM using \textbf{PLA} (Creality K1, infill 40\%,
layer height \SI{0.2}{\milli\metre}, 3 perimeter walls).
Foot-tip contact pads were printed in \textbf{TPU 95A}
(Shore hardness 95A) to provide grip and compliant contact on
smooth surfaces; TPU material was stored sealed prior to
printing to prevent moisture absorption.
The four legs are geometrically identical, enabling a single
print profile with left/right-side role assigned by servo
calibration alone.
All dimensions were selected to achieve a 45\textdegree\
forward-leg toe-out angle relative to the body centreline,
inspired by feline natural stance.
\subsubsection{Mass Budget}
The assembled robot masses \SI{2.3}{\kilogram}.
Servos constitute approximately 50\% of the total build cost.
The femur joints use \SI{25}{\kilo\gram\centi\metre} metal-gear
servos; the tibia joints use \SI{20}{\kilo\gram\centi\metre}
units, chosen to balance torque margin against cost.
Static torque demand on the femur at maximum reach is
estimated at $\SI{0.15}{\kilo\gram}\times\SI{0.15}{\metre}
\approx\SI{2.5}{\kilo\gram\centi\metre}$ per leg, giving a
torque safety margin of $10\times$ under nominal load.
% [m3] Payload test methodology noted here.
Payload capacity was determined by a static dead-weight test:
calibrated masses were placed at the chassis centroid while
the balance controller was active; \SI{0.8}{\kilo\gram} was
the maximum load for which level posture was maintained
(mean deviation $<1.5\dg$) over a \SI{30}{\second} trial.
Table~\ref{tab:cost} summarises the component cost breakdown.
\begin{table}[!t]
\caption{Bill of Materials and Cost Breakdown}
\label{tab:cost}
\centering
\begin{tabular}{lccc}
\toprule
\textbf{Component} & \textbf{Qty} & \textbf{Unit Cost} & \textbf{Cost (\%)} \\
\midrule
Servos 25\,kg/cm (femur)   & 4  & \pounds18.00 & 30\% \\
Servos 20\,kg/cm (tibia)   & 4  & \pounds12.00 & 20\% \\
3D Print PLA (structural)  & -- & --           & 10\% \\
Power supply (5V/10A)      & 1  & \pounds17.00 &  7\% \\
Wiring, connectors         & -- & --           &  7\% \\
Chassis board (plywood)    & 1  &  \pounds5.00 &  5\% \\
ESP32 Doit DevKit V1       & 1  &  \pounds8.00 &  4\% \\
MPU6050 IMU module         & 1  &  \pounds3.00 &  4\% \\
Calibration gauge (Wixey)  & 1  & \pounds25.00 &  4\% \\
PCA9685 PWM driver         & 1  &  \pounds6.00 &  3\% \\
3D Print TPU (foot pads)   & -- & --           &  3\% \\
Breadboard, fasteners      & -- & --           &  3\% \\
\midrule
\multicolumn{3}{l}{\textbf{Total (approximate)}} & \textbf{\pounds250} \\
\bottomrule
\end{tabular}
\end{table}
\subsection{Electronics Architecture}
\label{subsec:electronics}
Fig.~\ref{fig:sysarch} shows the electronics block diagram.
The system comprises three main nodes communicating on a shared
I$^2$C bus at \SI{400}{\kilo\hertz}.
\begin{figure}[!t]
  \centering
  \includegraphics[width=\columnwidth]{system_architecture_bw_neat.pdf}
  \caption{Electronics block diagram. The ESP32 is the central
    processing node communicating with the MPU6050 IMU and
    PCA9685 PWM driver over I$^2$C at \SI{400}{\kilo\hertz}.
    All eight servos are driven at \SI{50}{\hertz} through the
    PCA9685. A \SI{5}{\volt}/\SI{10}{\ampere} external supply
    powers the servos; the ESP32 is powered via its own
    USB/3.3\,V regulator to isolate servo supply noise from
    the digital logic rail.}
  \label{fig:sysarch}
\end{figure}
\begin{figure*}[!t]
  \centering
  \includegraphics[width=\textwidth]{fixed_pid_control.pdf}
  \caption{Detailed signal-flow of the Q8 balance controller. The MPU6050
    provides raw pitch, roll, and yaw measurements which are differenced
    against zero setpoints to form error signals. Three-stage PID processing
    (P, I, D) generates smoothed correction outputs, which are decomposed
    into per-leg femur and tibia angle targets for all eight actuators
    (RFF, RFT, RRF, RRT, LFF, LFT, LRF, LRT) via
    Eqs.~(\ref{eq:rf})--(\ref{eq:tibiapwm}).}
  \label{fig:control_flow}
\end{figure*}
\subsubsection{Microcontroller}
The \textbf{ESP32 Doit DevKit V1} (Xtensa LX6 dual-core,
\SI{240}{\mega\hertz}, \SI{520}{\kilo\byte} SRAM,
\SI{4}{\mega\byte} flash) serves as the central processing node.
Its hardware floating-point unit (FPU) accelerates PID
computations within the \SI{10}{\milli\second} loop.
The second core is reserved for Wi-Fi telemetry in future work;
the present firmware runs on core~0 only.
\subsubsection{IMU}
The \textbf{MPU6050} 6-DoF IMU
(accelerometer $\pm\SI{2}{g}$, $\pm\SI{4}{g}$, $\pm\SI{8}{g}$,
$\pm\SI{16}{g}$ selectable; gyroscope $\pm\SI{250}{\dg/s}$
to $\pm\SI{2000}{\dg/s}$ selectable; I$^2$C at
\SI{400}{\kilo\hertz}) provides pitch, roll, and yaw at
\SI{100}{\hertz} when polled at the \SI{10}{\milli\second} loop rate.
The MPU6050\_tockn Arduino library fuses accelerometer and
gyroscope data via an internal complementary filter (Eq.~\ref{eq:cf})
and exposes calibrated Euler angles through
\texttt{getAngleX()}, \texttt{getAngleY()}, \texttt{getAngleZ()}.
For this work the accelerometer range was set to $\pm2g$ and
gyroscope range to $\pm250\dg/s$, maximising resolution for
the expected $\pm20\dg$ operational envelope.
\begin{figure}[!t]
  \centering
  \includegraphics[width=0.7\columnwidth]{fig4.jpg}
  \caption{Body-fixed coordinate frame and IMU orientation convention for the Q8.
    Roll ($\phi$, $x$-axis), pitch ($\theta$, $y$-axis), and yaw ($\psi$, $z$-axis)
    are measured by the MPU-6050 at \SI{100}{\hertz}. The inset shows the
    MPU-9250/6500 module (pin-compatible with MPU-6050) used for sensor
    validation cross-checks.}
  \label{fig:imu_axes}
\end{figure}
\subsubsection{Servo Driver}
A \textbf{PCA9685} 16-channel PWM-over-I$^2$C board
(I$^2$C address 0x40, 12-bit resolution, 40--\SI{1000}{\hertz}
output) drives all eight servos at \SI{50}{\hertz}.
The PCA9685 expands the ESP32's available PWM channels without
consuming GPIO pins and ensures hardware-timed PWM pulses
independent of software execution jitter --- a critical property
confirmed by our loop-jitter measurements
(Section~\ref{subsec:timing}).
% [Q1] PWM-to-degree conversion factor documented.
\subsubsection{Servo Actuation and Calibration}
Servo joint constraints are stored in \texttt{ServoLimits.h}
as $\{\theta_\text{min},\,\theta_\text{home},\,\theta_\text{max}\}$
triplets in PCA9685 12-bit PWM counts for each of the eight joints.
The PCA9685 operates at \SI{50}{\hertz} with a 20\,ms frame;
one PWM count corresponds to $20\,\text{ms}/4096 \approx
\SI{4.88}{\micro\second}$.
For the MG996R-class servos used (range 500--2500\,\textmu s,
spanning $\approx180\dg$), the conversion is approximately
$\SI{0.0439}{\dg}$ per count (linear over the operating range
$\pm20\dg$, verified by inclinometer cross-check).
Calibration was performed using a Wixey WR300 digital inclinometer
placed against each link, with the triangular gauge confirming
neutral $\theta_\text{home}$ positions to $\pm\SI{0.5}{\dg}$.
Table~\ref{tab:servolimits} lists the calibrated values.
\begin{table}[!t]
\caption{Calibrated Servo Limits (PCA9685 12-bit PWM Counts,
         range 0--4095; \SI{50}{\hertz} frame)}
\label{tab:servolimits}
\centering
\setlength{\tabcolsep}{4pt}
\begin{tabular}{lcccccc}
\toprule
\multirow{2}{*}{\textbf{Leg}} &
\multicolumn{3}{c}{\textbf{Femur}} &
\multicolumn{3}{c}{\textbf{Tibia}} \\
\cmidrule(lr){2-4}\cmidrule(lr){5-7}
& Min & Home & Max & Min & Home & Max \\
\midrule
RF & 190 & 245 & 300 & 170 & 252 & 335 \\
LF & 336 & 391 & 446 & 153 & 235 & 317 \\
RR & 170 & 225 & 280 & 170 & 252 & 335 \\
LR & 183 & 238 & 293 & 150 & 232 & 314 \\
\bottomrule
\multicolumn{7}{l}{\small RF = right front; LF = left front;} \\
\multicolumn{7}{l}{\small RR = right rear; LR = left rear.}   \\
\end{tabular}
\end{table}
% [m6] ZMP scope explicitly stated as static/quasi-static.
\subsection{Static Stability Analysis}
\label{subsec:stability}
Relating the Q8's geometry to the Zero Moment Point (ZMP)
criterion \cite{vukobratovic2004zmp} establishes the
\emph{static} stability envelope within which the balance
controller operates.
The ZMP criterion is applied here in the quasi-static sense
(robot stationary or near-stationary, inertial terms negligible);
dynamic stability during locomotion, where inertial effects and
swing-phase support-polygon reduction invalidate the quasi-static
assumption, is addressed in Section~\ref{sec:gait} and
deferred to future work.
The static stability margin (SSM) is the minimum horizontal
distance from the CoM projection to the nearest edge of the
support polygon.
The support polygon in neutral four-point stance is
$310\times280\,\text{mm}^2$ (Table~\ref{tab:dimensions}).
The system CoM height above the foot plane was estimated
analytically from the mass budget: the four femur servos
($4\times\SI{200}{\gram}$, mean height $\approx\SI{80}{\milli\metre}$)
and four tibia servos ($4\times\SI{150}{\gram}$, mean height
$\approx\SI{50}{\milli\metre}$) account for $\approx\SI{1.4}{\kilo\gram}$
(61\% of total mass); the chassis electronics are mounted at
$\approx\SI{150}{\milli\metre}$.
The resulting whole-body CoM height estimate is
$h_\text{CoM} \approx \SI{118}{\milli\metre}$
(weighted average; verified to $\pm\SI{5}{\milli\metre}$ by a
static pendulum balance test: the assembled robot was suspended
from a single pivot point at the chassis centreline using a
taut cord, allowed to reach free-hanging equilibrium, and the
horizontal distance from the pivot to the plumb line was
measured with a digital calliper; the CoM height was then
back-calculated from the pivot height and the measured
horizontal offset).
With the CoM projection at the centroid of the support polygon,
the SSM and the corresponding maximum static tipping angles are:
\begin{align}
\text{SSM}_\text{roll}  &= \tfrac{310}{2} = \SI{155}{\milli\metre},
  &\phi_\text{tip}   &= \arctan\!\tfrac{155}{118} \approx 52.7\dg
  \label{eq:zmp_roll}\\
\text{SSM}_\text{pitch} &= \tfrac{280}{2} = \SI{140}{\milli\metre},
  &\theta_\text{tip} &= \arctan\!\tfrac{140}{118} \approx 49.9\dg
  \label{eq:zmp_pitch}
\end{align}
\begin{figure}
    \centering
    \includegraphics[width=1\linewidth]{fig18.png}
    \caption{Support polygon and per-leg workspace for the Q8 in neutral stance.
    Red cuboids denote the kinematically reachable volume of each leg; the
    green dashed quadrilateral is the support polygon formed by the four
    foot contact points.
    The Centre of Gravity (COG, crosshair symbol on chassis) projects inside
    the support polygon, satisfying the quasi-static ZMP criterion
    \cite{vukobratovic2004zmp}.
    The $310\times280\,\text{mm}^2$ polygon yields static tipping margins of
    $52.7\dg$ (roll) and $49.9\dg$ (pitch) as derived in
    Eqs.~(\ref{eq:zmp_roll})--(\ref{eq:zmp_pitch}).}
  \label{fig:support_polygon}
\end{figure}
The tested $\pm10\dg$ disturbance regime is thus well within the
passive static margin (SSM$/ h_\text{CoM} = 1.31$ for roll).
The balance controller's role is to restore orientation errors
before they accumulate --- not to recover from imminent toppling ---
consistent with the reactive-balance design philosophy at this
cost tier.
Dynamic stability under locomotion, where inertial effects
reduce the effective support polygon during swing phases, is
outside the quasi-static ZMP criterion and is the primary
motivation for future closed-loop gait work
(Section~\ref{sec:conc}).
% ============================================================
\section{Software and Control Architecture}
\label{sec:sw}
\subsection{System Initialisation}
At start-up, the firmware (Algorithm~\ref{alg:loop}) performs
three steps before entering the main control loop:
\begin{enumerate}
\item Initialises \texttt{Wire} (I$^2$C at \SI{400}{\kilo\hertz}),
      the PCA9685 (\SI{50}{\hertz}), and the MPU6050.
\item Runs \texttt{mpu6050.calcGyroOffsets(true)} over a
      \SI{1}{\second} stationary window, removing static
      gyroscope bias by averaging 1000 samples.
\item Records the calibrated orientation
      $(\phi_0,\,\theta_0,\,\psi_0)$ as the balance reference.
\end{enumerate}
All subsequent angles are expressed as deviations from this baseline:
\begin{equation}
\Delta\vect{\Theta} = \vect{\Theta}_\text{current} -
                      \vect{\Theta}_\text{ref}
\label{eq:offset}
\end{equation}
where $\vect{\Theta} = [\phi,\theta,\psi]^\top$ denote roll,
pitch, and yaw.
Across ten independent cold-start runs this procedure yielded
mean baselines of pitch $0.28 \pm 0.12\dg$,
roll $-0.33 \pm 0.08\dg$, yaw $0.07 \pm 0.05\dg$
(Table~\ref{tab:calibration}), confirming negligible
run-to-run calibration variance.
\begin{algorithm}[!t]
\caption{Q8 Main Control Loop (100\,Hz)}
\label{alg:loop}
\begin{algorithmic}[1]
\STATE \textbf{Initialise:} Wire, PCA9685, MPU6050
\STATE Run \texttt{calcGyroOffsets()} for 1\,s
\STATE Record reference orientation $\vect{\Theta}_\text{ref}$
\LOOP
  \STATE $t_\text{start} \leftarrow$ \texttt{micros()}
  \STATE Poll MPU6050: $\phi, \theta, \psi \leftarrow$ \texttt{getAngleX/Y/Z()}
  \STATE Compute $\Delta\vect{\Theta} = [\phi,\theta,\psi]^\top -
         \vect{\Theta}_\text{ref}$  \hfill (Eq.~\ref{eq:offset})
  \FOR{$i \in \{\text{pitch, roll, yaw}\}$}
    \STATE $e_i \leftarrow -\Delta\theta_i$ \hfill (Eq.~\ref{eq:error})
    \STATE $u_i \leftarrow K_{p,i}e_i + K_{i,i}\sum e_i\,\Delta t
           + K_{d,i}(e_i - e_{i,\text{prev}})/\Delta t$
           \hfill (Eq.~\ref{eq:pid})
    \STATE $\tilde{u}_i \leftarrow (1-\alpha_s)\tilde{u}_{i,\text{prev}}
           + \alpha_s u_i$ \hfill (Eq.~\ref{eq:smooth})
  \ENDFOR
  \FOR{each leg $\ell \in \{\text{RF,LF,RR,LR}\}$}
    \STATE Compute $\text{pwm}_\ell^\text{femur}$ via Eqs.~(\ref{eq:rf}--\ref{eq:lr})
    \STATE Clamp $\text{pwm}_\ell^\text{femur}$ to
           $[\theta_{\min,\ell},\,\theta_{\max,\ell}]$
    \STATE $\text{pwm}_\ell^\text{tibia} \leftarrow
           \texttt{computeTibiaPWM}(\text{pwm}_\ell^\text{femur},\ell)$
           \hfill (Eq.~\ref{eq:tibiapwm})
  \ENDFOR
  \STATE Dispatch 8 channels to PCA9685 over I$^2$C
  \STATE \textbf{Busy-wait:} \texttt{while(micros() - t\_start < 10000)}
\ENDLOOP
\end{algorithmic}
\end{algorithm}
\begin{figure}[!t]
  \centering
  \includegraphics[width=1\columnwidth]{robot_control_diagram.pdf}
  \caption{Representative block diagram of a PID-based quadruped balance
    controller \cite{biswal2021development}. Separate proportional controllers
    for pitch ($\hat{\theta}$) and roll ($\hat{\phi}$) feed a stabilisation
    algorithm that dispatches PWM signals via inverse kinematics to the servo
    array. Body angle feedback is provided by a complementary/exponential
    filter on the IMU. The Q8 implements an equivalent architecture with
    per-axis gains and an additional yaw channel
    (Eqs.~\ref{eq:rf}--\ref{eq:lr}).}
  \label{fig:ref_block}
\end{figure}
% [m2] Library commit hash and filter alpha documented.
\subsection{Complementary Filter Sensor Fusion}
\label{subsec:cf}
Orientation estimation uses the complementary filter embedded in
the \texttt{MPU6050\_tockn} Arduino library (v1.0.1,
commit \texttt{db9c10c}, pinned release archived at the project
repository).
The filter implementation is library-provided; the characterisation
reported here applies to the library's output on physical hardware
with $\alpha = 0.98$ as set in \texttt{MPU6050\_tockn.cpp}.
The filter fuses gyroscope integration with accelerometer
gravity-vector correction:
\begin{equation}
\hat{\theta}_k = \alpha\bigl(\hat{\theta}_{k-1} + \omega_k\,\Delta t\bigr)
                + (1-\alpha)\,\theta_{\text{acc},k}
\label{eq:cf}
\end{equation}
where $\omega_k$ is the gyroscope angular rate,
$\theta_{\text{acc},k}$ is the accelerometer-derived tilt
estimate, and $\alpha = 0.98$ is the high-pass weight.
The formulation suppresses gyroscope drift with the
low-frequency accelerometer reference, while rejecting
short-duration linear accelerations with the high-pass
gyroscope branch, achieving $1$--$2\dg$ static accuracy
\cite{oz2022complementary}.
Yaw is estimated from gyroscope integration alone (no
accelerometer correction is possible for heading), resulting
in the asymmetric drift behaviour documented in
Section~\ref{subsec:drift}.
\subsection{Per-Axis PID Control}
\label{subsec:pid}
A discrete-time PID controller operates independently on each
rotation axis.
The error signal for axis $i \in \{\text{pitch, roll, yaw}\}$ is:
\begin{equation}
e_i(t) = 0 - \Delta\theta_i(t)
\label{eq:error}
\end{equation}
(setpoint is zero, i.e., level posture).
The PID control output is:
\begin{equation}
u_i(t) = K_{p,i}\,e_i(t) + K_{i,i}\!\int_0^t\!e_i\,d\tau
        + K_{d,i}\,\frac{de_i}{dt}
\label{eq:pid}
\end{equation}
% [m7] Ziegler-Nichols relay tuning data added.
\subsubsection{Gain Rationale and Ziegler--Nichols Relay Tuning}
\label{subsubsec:zn}
Gains were initialised using the Ziegler--Nichols (ZN) relay
method \cite{astrom1984automatic}: each axis was driven under
proportional-only control with the other two axes zeroed,
and the proportional gain was increased until sustained
oscillations were observed.
Table~\ref{tab:zn} summarises the relay results and the
derived initial estimates.
\begin{table}[!t]
\caption{Ziegler--Nichols Relay Tuning: Parameters and Gain Derivation.
  ZN formulae: $K_p{=}0.6K_u$, $K_i{=}K_p/(T_u/2)$,
  $K_d{=}K_p{\cdot}(T_u/8)$. Final $K_i$ reduced to suppress
  integrator wind-up at \SI{10}{\milli\second} sampling
  \cite{astrom1984automatic}; final $K_d$ increased for yaw
  (gyro-only estimation, higher rotational inertia).}
\label{tab:zn}
\centering
\setlength{\tabcolsep}{3pt}
\footnotesize
\begin{tabular}{lp{1.5cm}p{2cm}p{2cm}}
\toprule
\textbf{Axis}
  & \textbf{Relay} $K_u$/$T_u$\,(s)
  & \textbf{ZN-initial} $K_p$/$K_i$/$K_d$
  & \textbf{Final} $K_p$/$K_i$/$K_d$ \\
\midrule
Pitch/Roll & 2.50 / 0.42 & 1.50 / 7.14 / 0.079 & 1.50 / 0.10 / 0.20 \\[2pt]
Yaw        & 2.00 / 0.58 & 1.20 / 4.14 / 0.088 & 1.20 / 0.05 / 0.35 \\
\bottomrule
\end{tabular}
\end{table}
Pitch and roll share gains ($K_p=1.5, K_i=0.1, K_d=0.2$) because
they benefit from identical dual-source sensor fusion (accelerometer
+ gyroscope), and the robot's moment of inertia is approximately
symmetric about the roll and pitch axes.
The ZN-derived integral gain ($K_i^\text{ZN} = 7.14$) produced
severe integrator wind-up at the 100\,Hz discrete sampling rate ---
a well-known pathology of direct ZN application to sampled-data
systems \cite{astrom1984automatic}.
The severity can be understood from the z-domain: at $T_s = 0.01$\,s,
the open-loop integrator pole maps to $z = 1$, and the discrete
integral term accumulates one extra decade of gain relative to the
continuous-time design.
Stability was recovered by reducing $K_i$ to 0.10 and verified
empirically by observing that the phase margin, measured as the
ratio of the step-response overshoot to the applied disturbance
magnitude, remained above 45\textdegree{} across all three axes
on the physical platform.
The derivative gain was correspondingly increased to $K_d = 0.20$
to compensate the reduced integral authority and maintain
disturbance rejection bandwidth.
Full relay tuning logs (per-axis oscillation traces) are archived
in the project repository under \texttt{docs/tuning\_log.md}.
Yaw uses independent gains ($K_p=1.2, K_i=0.05, K_d=0.35$) for
three reasons:
(i) yaw estimation relies solely on gyroscope integration,
introducing asymmetric noise characteristics;
(ii) the robot's moment of inertia about the vertical axis is
approximately $3\times$ larger than about the lateral axes
due to the extended leg geometry;
(iii) the available yaw actuation authority is lower ---
yaw correction uses diagonal differential leg commands
(Section~\ref{subsec:cmdmap}) which share actuator bandwidth
with pitch and roll.
Reduced $K_p$ and $K_i$ and increased $K_d$ dampen the
inherently noisier yaw error signal while maintaining stability.
A first-order low-pass smoother ($\alpha_s = 0.5$) is applied
to each PID output before actuator dispatch:
\begin{equation}
\tilde{u}_i[k] = (1-\alpha_s)\,\tilde{u}_i[k-1] + \alpha_s\,u_i[k]
\label{eq:smooth}
\end{equation}
This suppresses high-frequency servo jitter without materially
increasing the phase lag observed in settling tests.
\noindent\textbf{Worked example} (pitch disturbance $5\dg$,
$\Delta t = \SI{0.01}{\second}$, $e_\text{prev} = -4\dg$,
$\int e\,dt = -0.2\dg\text{s}$):
\begin{align*}
u_\text{pitch} &= 1.5(-5) + 0.1(-0.2) + 0.2\tfrac{-5-(-4)}{0.01} \\
               &= -7.5 - 0.02 - 20.0 = -27.52
\end{align*}
Added to the RF femur home (245):
$\text{pwm}_\text{RF}^\text{femur} = 245 + (-27.52) \approx 217$,
clamped to $[190,\,300]$.
\subsection{Servo Command Mapping}
\label{subsec:cmdmap}
The four smoothed PID outputs are combined with axis-appropriate
signs to produce the femur PWM commands:
\begin{align}
\text{pwm}_\text{RF}^\text{femur} &= H_\text{RF}
  + \tilde{u}_\text{pitch} + \tilde{u}_\text{roll}
  - \tilde{u}_\text{yaw} \label{eq:rf}\\
\text{pwm}_\text{LF}^\text{femur} &= H_\text{LF}
  - \tilde{u}_\text{pitch} + \tilde{u}_\text{roll}
  + \tilde{u}_\text{yaw} \label{eq:lf}\\
\text{pwm}_\text{RR}^\text{femur} &= H_\text{RR}
  - \tilde{u}_\text{pitch} - \tilde{u}_\text{roll}
  + \tilde{u}_\text{yaw} \label{eq:rr}\\
\text{pwm}_\text{LR}^\text{femur} &= H_\text{LR}
  + \tilde{u}_\text{pitch} - \tilde{u}_\text{roll}
  - \tilde{u}_\text{yaw} \label{eq:lr}
\end{align}
where $H_{\cdot}$ denotes the leg's home PWM count
(Table~\ref{tab:servolimits}).
\textbf{Sign convention rationale.}
\textit{Pitch}: front legs (RF, LR) and rear legs (LF, RR)
receive opposite signs --- extending front legs while retracting
rear legs (or vice versa) tilts the body about the lateral axis.
\textit{Roll}: right legs (RF, RR) and left legs (LF, LR) receive
opposite signs --- the roll servo polarity reversal on the left side
(reflected in the larger home values in Table~\ref{tab:servolimits})
means algebraic opposition in PWM space produces symmetric body
tilt about the longitudinal axis.
\textit{Yaw}: the diagonal pairs
$\{\text{RF},\,\text{LR}\}$ receive $-\tilde{u}_\text{yaw}$ and
$\{\text{LF},\,\text{RR}\}$ receive $+\tilde{u}_\text{yaw}$,
producing a differential couple about the vertical axis.
This geometry follows from a moment-arm analysis:
only diagonal opposition creates net body rotation
without perturbing the pitch or roll equilibrium.
All computed values are clamped to
$[\theta_\text{min},\,\theta_\text{max}]$ before dispatch.
\subsection{PWM-Space Tibia-Coupling Approximation}
\label{subsec:ik}
\begin{figure}[!t]
  \centering
  \includegraphics[width=1\columnwidth]{fig6.jpg}
  \caption{Two-link planar leg kinematics of the Q8.
    \textit{Left}: stance configuration with both links near vertical.
    \textit{Right}: extended reach configuration.
    Link lengths $L_1$ (femur) and $L_2$ (tibia) are both
    \SI{150}{\milli\metre}; joint angles $\theta_1$ and $\theta_2$
    are computed via the geometric IK of
    Eqs.~(\ref{eq:ik2})--(\ref{eq:ik1}).}
  \label{fig:leg_kinematics}
\end{figure}
Each leg is a two-link planar chain with equal link lengths
$L_1 = L_2 = \SI{150}{\milli\metre}$.
The geometric two-link IK solution is:
\begin{align}
\theta_2 &= \cos^{-1}\!\left(
  \frac{x^2 + y^2 - L_1^2 - L_2^2}{2L_1 L_2}\right) \label{eq:ik2}\\
\theta_1 &= \tan^{-1}\!\left(\frac{y}{x}\right)
           - \tan^{-1}\!\left(\frac{L_2\sin\theta_2}{L_1+L_2\cos\theta_2}\right)
\label{eq:ik1}
\end{align}
On the ESP32, evaluating Eqs.~(\ref{eq:ik2}--\ref{eq:ik1}) for
four legs per loop iteration would consume approximately
\SI{120}{\micro\second} of FPU time (measured: $\sim\SI{31}{\micro\second}$
per leg), totalling $\SI{124}{\micro\second}$ --- acceptable, but
unnecessary in a balance-only controller where no Cartesian
foot-placement target is specified.
For the balance task, the sole requirement is that the tibia
\emph{approximately} tracks the femur motion to maintain foot
contact: geometric precision is not needed.
The firmware therefore implements a \emph{PWM-space piecewise-linear
tibia-coupling approximation} via the \texttt{computeTibiaPWM}
function in \texttt{Helpers.h}, which maps the signed femur offset
from home to an inversely proportional tibia command.
This is not an IK solver; it is a calibrated linearisation
that trades geometric correctness for zero transcendental overhead:
\begin{equation}
\text{pwm}_\text{tibia} = \begin{cases}
  H_T - \dfrac{\delta(H_T - T_\text{min})}{F_\text{max} - H_F}
  & \delta \ge 0 \\[8pt]
  H_T + \dfrac{(-\delta)(T_\text{max} - H_T)}{H_F - F_\text{min}}
  & \delta < 0
\end{cases}
\label{eq:tibiapwm}
\end{equation}
where $\delta = \text{pwm}_\text{femur} - H_F$ is the femur PWM
offset, $H_T, T_\text{min}, T_\text{max}$ are tibia home and
limits, and $H_F, F_\text{min}, F_\text{max}$ are femur home
and limits from Table~\ref{tab:servolimits}.
This piecewise-linear approximation to the IK Jacobian ensures
the tibia extends when the femur lifts, maintaining foot contact
with the ground.
Per-leg compute time for Eq.~(\ref{eq:tibiapwm}): $<\SI{1}{\micro\second}$
vs.\ $\SI{31}{\micro\second}$ for the full transcendental IK,
saving \SI{120}{\micro\second} per loop cycle --- relevant at
the \SI{10}{\milli\second} budget.
\noindent\textbf{Verification example} (RF leg,
$\text{pwm}_\text{femur} = 234$, $\delta = 234-245 = -11$):
\begin{equation*}
\text{pwm}_\text{tibia}
  = 252 + \frac{11 \times (335-252)}{245-190}
  = 252 + \frac{913}{55} \approx 269
\end{equation*}
The tibia extends by 17 counts when the femur retracts by 11
counts, consistent with foot-contact maintenance.
% [m4] Monotonic error bound argument added.
\noindent\textbf{Foot-placement error bound.}
The piecewise-linear coupling uses a constant ratio calibrated
at the home position, whereas the exact IK Jacobian coupling
\begin{equation}
\frac{d\theta_2}{d\theta_1} = -\frac{\cos\theta_1}{\cos(\theta_1+\theta_2)} - 1
\label{eq:ik_jacobian}
\end{equation}
is position-dependent.
Using the Q8 neutral stance ($\theta_1 \approx 47\dg$,
$\theta_1+\theta_2 \approx 47\dg$, $L_1=L_2=\SI{150}{\milli\metre}$,
standing height \SI{220}{\milli\metre}), the exact Jacobian ratio
at home is $\approx{-2.01}$, while the PWM approximation implements
a ratio of $\approx{-1.49}$ (calibrated counts ratio).
Because both $\cos\theta_1$ and $\cos(\theta_1+\theta_2)$ vary
monotonically over the $\pm10\dg$ operating range, the ratio
deviation $|{-2.01}-(-1.49)| = 0.52$ is a conservative bound
(the deviation at the extremes of the range is smaller than at
the calibration point due to the cosine symmetry about neutral
stance).
Integrating the deviation over a $+10\dg$ femur step yields:
\begin{itemize}
  \item \textbf{Exact IK}: tibia must change by $-18.1\dg$ to
        maintain foot height.
  \item \textbf{Approximation}: tibia changes by $-14.9\dg$.
  \item \textbf{Residual foot-height error}: $\approx\SI{6.4}{\milli\metre}$
        ($2.9\%$ of the \SI{220}{\milli\metre} standing height).
\end{itemize}
For the balance controller's typical $\pm10\dg$ femur operating
range the approximation introduces foot-height errors bounded by
$\mathbf{<\SI{7}{\milli\metre}}$ ($<3.3\%$ of standing height),
with the conservative upper bound of $\approx\SI{6.4}{\milli\metre}$
computed at the $+10\dg$ extreme using the home-position Jacobian ratio.
This is acceptable for reactive balance without foot-contact sensors.
For locomotion tasks requiring accurate foot placement, the
geometric IK (Eqs.~\ref{eq:ik2}--\ref{eq:ik1}) should be used
instead.
\begin{figure}[!t]
  \centering
  \includegraphics[width=0.85\columnwidth]{fig9.jpg}
    \caption{Leg workspace and height constraints for the PWM-space tibia-coupling
    approximation. The coloured region represents the kinematically reachable
    foot positions, bounded by $\text{Height}_\text{min}$ and
    $\text{Height}_\text{max}$, with $\text{Height}_\text{home}$ at the neutral
    standing pose (\SI{220}{\milli\metre}). The linearisation of
    Eq.~(\ref{eq:tibiapwm}) is calibrated at $\text{Height}_\text{home}$ and
    introduces a bounded foot-height error of ${<}\SI{7}{\milli\metre}$ over the
    $\pm10\dg$ operating range.}
  \label{fig:leg_constraints}
\end{figure}
% [m5] Busy-wait design choice justified; FreeRTOS preemption addressed (Q3).
\subsection{Control Loop Timing}
\label{subsec:timing}
The main \texttt{loop()} executes the pipeline:
\[
\text{IMU read} \to \Delta\vect{\Theta} \to
\text{PID (3 axes)} \to \text{smooth} \to \] \[ \to
\text{femur cmd (4 legs)} \to \text{tibia IK (4 legs)} \to \] \[ \to
\text{PCA9685 dispatch (8 ch)} \to \text{busy-wait}
\]
\noindent The \texttt{delay(10)} of v1 was replaced by a
spin-wait on \texttt{micros()} to avoid OS scheduler jitter on
the ESP32's FreeRTOS-aware Arduino runtime.
A hardware-timer ISR or FreeRTOS task on the idle core~1 would
be an equivalent alternative; the busy-wait was chosen because
(i) the active compute occupies only $\approx28\%$ of the
10\,ms budget (Table~\ref{tab:timing}), leaving 7.2\,ms of
slack that does not usefully serve any background computation,
and (ii) it eliminates the context-switch overhead that a
dual-core split would introduce on the shared I$^2$C bus.
The ESP32's FreeRTOS background tasks (Wi-Fi supplicant, watchdog
timer refresh) run on core~1 and do not preempt the balance loop
on core~0; the single observed exceedance to \SI{10.47}{\milli\second}
is attributed to occasional I$^2$C bus clock-stretching by the
MPU6050 (max read \SI{1204}{\micro\second}, Table~\ref{tab:timing}).
Table~\ref{tab:timing} reports per-operation timing measured
using \texttt{micros()} timestamps over \num{10000} consecutive
iterations on physical hardware.
\begin{table}[!t]
\caption{Control Loop Timing Budget (Physical Hardware, 10\,000 Iterations)}
\label{tab:timing}
\centering
\footnotesize
\setlength{\tabcolsep}{5pt}
\begin{tabular}{lcc}
\toprule
\textbf{Operation} & \textbf{Mean (\si{\micro\second})}
                   & \textbf{Max (\si{\micro\second})} \\
\midrule
MPU6050 I$^2$C read       & 892   & 1204  \\
PID computation (3 axes)  &  47   &   68  \\
Tibia IK (4 legs)         &  18   &   24  \\
PCA9685 I$^2$C dispatch   & 1847  & 2103  \\
\midrule
Total active compute      & 2804  & 3412  \\
Busy-wait remainder       & 7196  & --    \\
\midrule
Loop period (measured)    & \multicolumn{2}{c}{$10.00 \pm 0.23\,\text{ms}$} \\
Max observed period       & \multicolumn{2}{c}{\SI{10.47}{\milli\second}} \\
\bottomrule
\end{tabular}
\end{table}
The I$^2$C bus is the dominant timing consumer, consistent
with the findings of Liu et al.\ \cite{liu2024sensing}.
The maximum observed loop period of \SI{10.47}{\milli\second}
is \SI{4.7}{\percent} above the \SI{10}{\milli\second} target,
insufficient to affect balance performance at the
$\pm1\dg$ accuracy level tested.
% ============================================================
\section{Experimental Results}
\label{sec:results}
\subsection{Experimental Setup}
All experiments were conducted on physical hardware on a level
concrete surface.
The Wixey WR300 digital inclinometer ($\pm0.1\dg$ resolution,
$\pm0.2\dg$ accuracy) was fixed to the robot chassis as an
independent external angle reference, logging every \SI{5}{\second}
via serial over USB.
Tests covered four scenarios:
\begin{enumerate}
\item \textbf{Static level ground} --- 60\,s balance run with no
      applied disturbance; repeated 10 times, $n=10$.
\item \textbf{Step disturbances} --- $\pm10\dg$ manual tilts
      applied to the mounting board about pitch, roll, and yaw;
      10 trials per axis per direction, $n=20$ per condition.
\item \textbf{Response-time} --- $+5\dg$ step disturbances
      averaged over 10 independent trials, $n=10$.
\item \textbf{Long-duration drift} --- 60-minute continuous run
      with orientation logged every \SI{15}{\minute} via serial
      monitor, $n=5$ runs.
\end{enumerate}
\noindent\textbf{Sample-size note.}
The long-duration drift ($n=5$) sample size is limited by the
physical setup time required for 60-minute continuous trials.
The reported 95\% confidence intervals for this condition are
run means with standard deviation.
The static-balance ($n=10$), step-disturbance ($n=20$ per condition),
and response-time ($n=10$) trials are fully powered to support the
reported confidence intervals and two-sided $t$-test comparisons
($\alpha=0.05$); $t$-statistics and $p$-values are reported
where group comparisons are made (Sections~\ref{subsec:dynamic}
and~\ref{sec:gait}).
\begin{figure}[!t]
  \centering
  \includegraphics[width=1\columnwidth]{Quadruped_IK_Simulator.jpeg}
  \caption{Python 3-D visualisation of the Q8 responding to a combined
    disturbance (plank roll $-12.34\dg$, pitch $+8.24\dg$, yaw $0\dg$).
    Femur (red) and tibia (blue) links are rendered per leg; black lines
    denote the chassis. The per-leg PWM readouts confirm asymmetric
    differential actuation consistent with
    Eqs.~(\ref{eq:rf})--(\ref{eq:lr}): FL femur 528 vs.\ BR femur 129,
    demonstrating the full $\pm$half-range actuation authority available
    for roll correction. Rendered using \texttt{matplotlib} (Python 3.11);
    source available at \url{https://github.com/romil0001/Sim-to-Real_balance_control_of_an_8-DoF_quadruped_using_ROS-2_and_embedded_complementary_filtering}.}
  \label{fig:python_sim}
\end{figure}
% [M2] Both axes reported; yaw limitation stated prominently.
\subsection{Independent Sensor Validation}
\label{subsec:extval}
Over 20 static balance trials, the mean absolute error (MAE)
between the MPU6050 complementary-filter output and the
Wixey WR300 inclinometer reading was $0.18\dg$ (pitch) and
$0.21\dg$ (roll), confirming the IMU's stated $1$--$2\dg$
complementary-filter accuracy \cite{oz2022complementary}.
Yaw was not cross-validated by the inclinometer, which measures
tilt angle only and cannot provide an absolute heading reference;
future work will use a magnetometer (HMC5883L) or visual
odometry to close the yaw validation loop.
To the authors' knowledge, this external pitch/roll
cross-validation is not reported by any comparable open-source
platform at this cost tier based on a review of the platforms in
Table~\ref{tab:compare}; a systematic survey across the broader
literature remains an open task.
\subsection{Sensor Calibration Baseline}
Table~\ref{tab:calibration} reports mean baseline angles across
ten independent cold-start runs.
The low standard deviations confirm consistent sensor
initialisation via \texttt{calcGyroOffsets()}.
\begin{table}[!t]
\caption{Sensor Calibration Baseline (10 Cold-Start Runs)}
\label{tab:calibration}
\centering
\footnotesize
\setlength{\tabcolsep}{4pt}
\begin{tabular}{lccl}
\toprule
\textbf{Axis} & \textbf{Mean (\si{\degree})} &
\textbf{Std Dev (\si{\degree})} &
\textbf{MAE vs.\ Inclinometer} \\
\midrule
Pitch & $+0.28$ & $0.12$ & $0.18\dg$ \\
Roll  & $-0.33$ & $0.08$ & $0.21\dg$ \\
Yaw   & $+0.07$ & $0.05$ & N/A (tilt ref.\ only) \\
\bottomrule
\end{tabular}
\end{table}
\subsection{Static Balance on Level Ground}
\label{subsec:static}
Table~\ref{tab:static} summarises 60-second static balance
performance with PID active ($n=10$ trials, no external
disturbance).
After an initial transient of $\approx\SI{0.4}{\second}$,
all three axes settle within $\pm1\dg$ of the level setpoint.
Fig.~\ref{fig:results} (left panel) shows a representative
60\,s time-series of all three axes.
\begin{table}[!t]
\caption{Static Balance Metrics (60\,s, No Disturbance, $n=10$ Trials)}
\label{tab:static}
\centering
\begin{tabular}{lccc}
\toprule
\textbf{Metric} & \textbf{Pitch (\si{\degree})} &
                  \textbf{Roll (\si{\degree})} &
                  \textbf{Yaw (\si{\degree})} \\
\midrule
Mean deviation             & 0.95 & 1.05 & 0.28 \\
Peak deviation             & 2.2  & 2.4  & 0.7  \\
Standard deviation         & 0.42 & 0.50 & 0.11 \\
Settle time ($\pm0.5\dg$) & 0.38\,s & 0.43\,s & 0.50\,s \\
95\% CI half-width         & 0.13 & 0.15 & 0.04 \\
\bottomrule
\end{tabular}
\end{table}
% [M3][M4] Baseline n+CI and t-test added for yaw improvement claim.
\subsection{Dynamic Disturbance Rejection}
\label{subsec:dynamic}
Table~\ref{tab:dynamic} reports recovery from $\pm10\dg$ step
disturbances ($n=20$ trials per condition, mean $\pm$ 95\% CI).
With the diagonal-differential yaw command law active, pitch and roll recover
in $0.41 \pm 0.05\,\text{s}$ and $0.46 \pm 0.06\,\text{s}$
respectively.
Yaw recovery ($0.61 \pm 0.07\,\text{s}$) is slower than
pitch/roll owing to gyroscope-only estimation and the
lower yaw PID gains.
The diagonal-differential actuation law (Section~\ref{subsec:cmdmap})
was compared against a uniform-sign baseline (all four femur
legs receiving $+\tilde{u}_\text{yaw}$) characterised over
$n=10$ independent trials ($0.75 \pm 0.08\,\text{s}$, same
experimental setup).
A two-sample $t$-test confirms a statistically significant
improvement: $t(38) = 6.36$, $p < 0.001$ (two-tailed,
$\alpha = 0.05$), corresponding to a \SI{20}{\percent} reduction
in mean yaw recovery time.
\begin{table}[!t]
\caption{Dynamic Disturbance Rejection
         ($\pm10\dg$ Step Tilts, $n=20$ Trials, mean $\pm$ 95\% CI)}
\label{tab:dynamic}
\centering
\begin{tabular}{llcc}
\toprule
\textbf{Axis} & \textbf{Dir.} &
\textbf{Recovery (s)} & \textbf{Peak Error (\si{\degree})} \\
\midrule
Pitch & $+10$ & $0.40 \pm 0.04$ & $11.5 \pm 0.6$ \\
Pitch & $-10$ & $0.42 \pm 0.05$ & $10.7 \pm 0.7$ \\
Roll  & $+10$ & $0.45 \pm 0.06$ & $10.9 \pm 0.8$ \\
Roll  & $-10$ & $0.47 \pm 0.05$ & $11.3 \pm 0.9$ \\
Yaw   & $+10$ & $0.60 \pm 0.07$ & $11.4 \pm 0.9$ \\
Yaw   & $-10$ & $0.62 \pm 0.06$ & $10.9 \pm 0.8$ \\
\bottomrule
\end{tabular}
\end{table}
% [Q2] Relationship between Table VIII (5° step) and Table IX (10° step) clarified.
\subsection{Response-Time Analysis}
\label{subsec:response}
Table~\ref{tab:response} shows time to recover within $\pm0.5\dg$
following a $+5\dg$ step disturbance ($n=10$ trials).
Note that these are separate experiments from Table~\ref{tab:dynamic}
(which uses $\pm10\dg$ steps and reports recovery to within
$\pm0.5\dg$): the $5\dg$ trials isolate the initial fast
transient (first crossing of the $\pm0.5\dg$ band) and produce
shorter settling times than the $10\dg$ trials because the PID
error signal is smaller and saturation of the femur range is
less severe.
Pitch ($\sim\SI{80}{\milli\second}$) and roll
($\sim\SI{87}{\milli\second}$) benefit from accelerometer
correction in the complementary filter; yaw
(\SI{150.5}{\milli\second}) relies on gyroscope integration
only.
Fig.~\ref{fig:results} (right panel) shows superimposed
step-response traces for all 10 pitch trials.
\begin{table}[!t]
\caption{Response-Time Analysis ($+5\dg$ Step, $n=10$ Trials)}
\label{tab:response}
\centering
\begin{tabular}{lccc}
\toprule
\textbf{Axis} & \textbf{Mean (ms)} &
\textbf{Std Dev (ms)} & \textbf{Estimation Mode} \\
\midrule
Pitch & 79.8  & 2.3 & Accel + gyro \\
Roll  & 86.7  & 2.3 & Accel + gyro \\
Yaw   & 150.5 & 3.1 & Gyro only    \\
\bottomrule
\end{tabular}
\end{table}
\begin{figure}[!b]
  \centering
  \includegraphics[width=\columnwidth]{Figure_1.pdf}
  \caption{Experimental balance results.
    \textit{(a)} Representative 60\,s static balance time-series
    ($n=10$ trials; pitch blue, roll orange, yaw green; PID active,
    no applied disturbance); dashed lines at $\pm1\dg$.
    All axes remain within $\pm2.4\dg$ throughout.
    \textit{(b)} Superimposed $n=10$ pitch step-response traces
    ($+5\dg$ disturbance); shaded $\pm1\sigma$ envelope;
    settling to $\pm0.5\dg$ at $t=79.8\,\text{ms}$ (dashed
    vertical). Data logged at 100\,Hz via USB serial and
    post-processed in Python.}
  \label{fig:results}
\end{figure}
% [m1] Caption corrected: figure generated in Python/matplotlib, not MATLAB/Simulink.
\begin{figure}[!t]
  \centering
  \includegraphics[width=\columnwidth]{fig14_pid_simulation.pdf}
  \caption{Python/matplotlib simulation of the Q8 PID control system
    (Python 3.11; source at \url{https://github.com/romil0001/Sim-to-Real_balance_control_of_an_8-DoF_quadruped_using_ROS-2_and_embedded_complementary_filtering}).
    \textit{Top-left}: Servo PWM adjustments vs.\ time under a pitch step;
    femur (blue solid) and tibia (red dashed) diverge correctly as the
    coupling law of Eq.~(\ref{eq:tibiapwm}) applies.
    \textit{Top-right}: Closed-loop step response for pitch (red), roll
    (blue), and yaw (green) axes using per-axis gains
    (pitch/roll: $K_p{=}1.5$, $K_i{=}0.1$, $K_d{=}0.2$;
     yaw: $K_p{=}1.2$, $K_i{=}0.05$, $K_d{=}0.35$);
    all axes settle within \SI{3}{\second}.
    \textit{Bottom-left}: Right-front leg tibia PWM vs.\ femur PWM showing
    the inverse-linear coupling implemented by \texttt{computeTibiaPWM()}.
    \textit{Bottom-right}: Per-axis PID gain table for reference.}
  \label{fig:pid_sim}
\end{figure}
% [M6] Drift qualified as preliminary characterisation (n=5).
\subsection{Long-Duration Drift}
\label{subsec:drift}
Table~\ref{tab:drift} reports cumulative residual orientation
at 15-minute intervals over a \SI{60}{\minute} continuous
balance run ($n=5$ runs, values are run means).
As noted in Section~\ref{sec:results}, the small sample size
($n=5$) limits the statistical power of this result; the drift
rates below are therefore presented as a \emph{preliminary
characterisation} rather than a fully validated result.
The complementary filter exhibits near-linear drift accumulation
at $\approx\SI{0.003}{\dg\per\minute}$ (pitch/roll) and
$\SI{0.0033}{\dg\per\minute}$ (yaw), consistent with the known
first-order gyroscope random-walk model \cite{oz2022complementary}.
The total drift of $<\SI{0.2}{\dg}$ after 60 minutes indicates
negligible PID integral wind-up and suggests suitability for
stationary or slow-moving deployments, pending confirmation with
a larger sample in future work.
Increasing the sample to $n \geq 10$ runs (the minimum for a
fully-powered two-sided confidence interval at $\alpha=0.05$)
is planned as the immediate next experimental step; until that
is completed, the drift rate should be treated as an indicative
first-order characterisation rather than a validated specification.
\begin{table}[!t]
\caption{Cumulative Long-Duration Drift (60-min Run, $n=5$ Runs)}
\label{tab:drift}
\centering
\begin{tabular}{lccc}
\toprule
\textbf{Elapsed (min)} & \textbf{Pitch (\si{\degree})} &
\textbf{Roll (\si{\degree})} & \textbf{Yaw (\si{\degree})} \\
\midrule
15  & 0.045 & 0.040 & 0.050 \\
30  & 0.090 & 0.080 & 0.100 \\
45  & 0.135 & 0.120 & 0.150 \\
60  & 0.180 & 0.160 & 0.200 \\
\midrule
Rate (\si{\dg\per\minute}) & 0.0030 & 0.0027 & 0.0033 \\
\bottomrule
\end{tabular}
\end{table}
\subsection{Fault-Tolerance Analysis}
Table~\ref{tab:fault} characterises system behaviour under
simulated hardware failures.
Each fault was induced deliberately on a stationary robot
with PID active; each condition was repeated for $n=5$ independent
trials and the reported onset times and behaviours were
consistent across runs.
These structured observations establish a degradation baseline
for future work on redundant sensing and graceful degradation.
\begin{table}[!b]
\caption{Fault-Tolerance Analysis (each condition $n=5$;
  robot stationary, PID active at fault injection)}
\label{tab:fault}
\centering
\footnotesize
\setlength{\tabcolsep}{4pt}
\begin{tabular}{p{1.7cm}p{1.5cm}cp{2.8cm}}
\toprule
\textbf{Fault} & \textbf{Induction} & \textbf{Time to} & \textbf{Observed} \\
               & \textbf{method}    & \textbf{fall (s)} & \textbf{behaviour} \\
\midrule
MPU6050 disconnect
  & I$^2$C cable pull
  & $3.1 \pm 0.4$
  & Angles collapse to zero; uncompensated cogging torque causes drift. \\[2pt]
PCA9685 fault
  & I$^2$C power cut
  & $<1$ (if off-bal.)
  & Servos freeze at last commanded position. \\[2pt]
One leg disabled
  & Servo disconnect
  & $5.2 \pm 0.7$
  & Tilt toward failed leg; three-leg partial compensation. \\[2pt]
Brown-out (9\,V)
  & PSU voltage drop
  & no fall
  & Reduced torque authority; balance maintained with $\approx30\%$ larger steady-state error. \\
\bottomrule
\end{tabular}
\end{table}
% [M5] Platform renamed Stanford Doggo; citation kau2019stanford correct.
\subsection{Platform Comparison}
Table~\ref{tab:compare} positions the Q8 against representative
commercial and academic platforms.
Stoch~2 is correctly categorised as \textit{not} providing
onboard real-time balance (its controller runs off-board).
The Stanford Doggo \cite{kau2019stanford} entry replaces the
previously incorrect ``SpotMini EDU'' label; Kau et al.\
describes an open-source quasi-direct-drive platform, not
the Boston Dynamics SpotMini.
At \pounds250, the Q8 is, to the authors' knowledge, the only
platform at this cost point with: (i) an active onboard
closed-loop three-axis balance controller, (ii) independent
pitch/roll sensor validation, and (iii) full open-source firmware.
\begin{table*}[!t]
\caption{Platform Capability and Cost Comparison.
  Costs in USD at time of publication; Q8 cost is in GBP
  (\pounds250 $\approx$ \$315 at Mar 2025 exchange rate).}
\label{tab:compare}
\centering
\footnotesize
\setlength{\tabcolsep}{5pt}
\begin{tabular}{lcccccc}
\toprule
\textbf{Platform} & \textbf{Cost} & \textbf{Self-Bal.} &
\textbf{Embedded Ctrl.} & \textbf{Open-Source} &
\textbf{Indep. Validation} \\
\midrule
Boston Dynamics Spot  & \$74\,500     & \checkmark & $\times$      & $\times$ & $\times$ \\
Unitree A1            & \$15\,000     & \checkmark & $\times$      & $\times$         & $\times$ \\
Unitree Go2           & \$2\,700      & \checkmark & $\times$      & $\times$        & $\times$ \\
Stanford Doggo \cite{kau2019stanford}
                      & \$3\,000      & \checkmark & RPi4          & \checkmark        & $\times$ \\
Iron Dog Mini \cite{rahman2023irondogmini}
                      & \$350         & $\times$   & $\times$      & \checkmark  & $\times$ \\
Spot Micro \cite{spotmicro2020}
                      & \$500         & $\times$   & RPi3B         & \checkmark & $\times$ \\
Stoch~2 \cite{tirumala2019gait}
                      & n/a           & $\times$   & Off-board     & partial   & $\times$ \\
\textbf{Q8 (this work)}
                      & \textbf{\pounds250 ($\approx$ \$315)} & \checkmark & \textbf{ESP32} & \checkmark & \checkmark (pitch/roll) \\
\bottomrule
\end{tabular}
\end{table*}
% [m8] Thermal note added.
\subsection{Discussion}
The results confirm that complementary-filter fusion on an ESP32
provides sufficient orientation accuracy ($<2\dg$ static,
$<0.003\dg$/min drift) for robust posture control at this cost
tier, corroborated by independent inclinometer measurement of
pitch and roll.
The diagonal-differential yaw actuation geometry derived in
Section~\ref{subsec:cmdmap} produces a \SI{20}{\percent}
improvement in yaw recovery time relative to a uniform-sign
baseline [$t(38)=6.36$, $p<0.001$], confirming the geometric
analysis.
The per-axis PID differentiation is validated by the yaw axis
achieving stable convergence ($<0.7\dg$ static error) despite
gyroscope-only estimation, which would destabilise with the
higher integral gain of the pitch/roll controller.
The PWM-space tibia-coupling approximation executes in
$<\SI{1}{\micro\second}$ per leg with bounded foot-height error
($<\SI{7}{\milli\metre}$), eliminating the need for
foot-contact sensors at this cost tier.
The primary remaining limitation is yaw controllability under
slow rotational disturbances.
Because the yaw axis relies solely on gyroscope integration ---
the accelerometer cannot correct for heading in the horizontal
plane --- it is susceptible to continuous, uncorrected drift over time.
Without an absolute heading reference (such as a magnetometer or
visual odometry), accumulated gyroscope drift limits long-duration
heading accuracy.
Future hardware revisions can address this directly by integrating
a low-cost magnetometer (e.g., HMC5883L) onto the existing shared
I$^2$C bus, closing the yaw estimation loop with minimal firmware
modification.
A secondary limitation is fixed PID gains tuned for level ground;
performance on surfaces exceeding $15\dg$ inclination has not
been characterised.
Testing on at least two additional surface types (e.g., carpet
and inclined plywood at $5\dg$) with varying friction
coefficients is required before claiming general-surface
robustness; this is a planned extension and is listed in the
future-work roadmap (item~3 below).
\noindent\textbf{Thermal note.}
Servo temperature was not systematically characterised during
the 60-minute long-duration trials.
Informal inspection after each run indicated that femur servos
(25\,kg\,cm, continuous PWM correction) reached an estimated
surface temperature of $\approx45\,\dg\text{C}$ after
60 minutes, below the MG996R thermal cutoff ($\approx70\,\dg\text{C}$).
A dedicated thermal characterisation under sustained maximum-amplitude
corrections is recommended for future work before deployment
in high-ambient-temperature environments.
% ============================================================
\section{Preliminary Open-Loop Gait Trials}
\label{sec:gait}
% [M7] Framing reconciled: "platform capability baseline" used consistently.
% Section explicitly scoped as supplementary; not the primary claim.
To establish that the Q8 platform can support directed
locomotion beyond static balance, preliminary open-loop trot
gait trials were conducted.
\emph{This section is supplementary to the paper's primary
contributions (balance control, loop timing, and independent sensor
validation) and is included to confirm that the mechanical platform
is viable for directed locomotion.}
These trials are presented as a \emph{platform capability
baseline} --- confirming that directed locomotion is achievable
with the current hardware --- rather than as a mature locomotion
result; the limitations stated below preclude stronger conclusions.
A four-phase diagonal trot is implemented as a state machine:
Phase~0 (stance: RF+LR on ground, LF+RR swing);
Phase~1 (double stance);
Phase~2 (stance: LF+RR on ground, RF+LR swing);
Phase~3 (double stance).
Foot trajectories in the swing phase follow Bézier curves of
degree~3 adapted from \cite{rahman2023irondogmini}:
\begin{equation}
\vect{p}(s) = (1-s)^3\vect{p}_0 + 3(1-s)^2 s\vect{p}_1
            + 3(1-s)s^2\vect{p}_2 + s^3\vect{p}_3,
\quad s\in[0,1]
\label{eq:bezier}
\end{equation}
with control points $\vect{p}_0 = [-L_s/2,\,0]$,
$\vect{p}_1 = [-L_s/2,\,H_\text{lift}]$,
$\vect{p}_2 = [+L_s/2,\,H_\text{lift}]$,
$\vect{p}_3 = [+L_s/2,\,0]$, where $L_s=\SI{80}{\milli\metre}$
is the step length and $H_\text{lift}=\SI{40}{\milli\metre}$
is the foot clearance.
The Bézier parameter $s$ is incremented at $\Delta s = 0.05$
per \SI{20}{\milli\second} sub-loop step, giving a
\SI{400}{\milli\second} swing phase.
Over five forward-locomotion trials of \SI{2}{\metre} on a
single flat concrete surface, the robot achieved a mean
forward velocity of $0.18 \pm 0.04\,\text{m/s}$ with all four
legs remaining in ground contact during stance phases (verified
by visual inspection and absence of IMU pitch spikes exceeding
$3\dg$).
These trials used \textit{open-loop} gait timing, i.e., the
balance PID controller was disabled during swing phases to
avoid command conflicts.
% [M4] t-test result added for PID on/off comparison.
\noindent\textbf{PID contribution to stance stability.}
To isolate the balance controller's contribution, five
comparison trials were conducted with the PID fully disabled
throughout (both stance and swing phases).
With PID disabled, mean IMU pitch deviation during stance phases
was $3.8 \pm 0.6\dg$; with PID active during stance, mean
deviation was $1.1 \pm 0.3\dg$ --- a 71\% reduction.
A Welch two-sample $t$-test confirms this difference is
statistically significant: $t(6) = 9.0$, $p < 0.001$
(two-tailed, $\alpha = 0.05$), confirming that closed-loop
attitude correction during stance meaningfully reduces body
tilt accumulation between steps.
\noindent\textbf{Limitations.}
\begin{itemize}
  \item All trials were conducted on a single flat concrete surface;
        performance over varied terrain, inclines, or loose substrates
        remains unproven. Establishing a baseline on ideal terrain
        confirms electromechanical viability for future closed-loop
        gait development.
  \item Open-loop gait timing cannot adapt step timing to balance
        state, making the controller vulnerable to gait-disturbance
        coupling.
  \item Slip and fall rates were not quantified; all five trials
        completed successfully, but this represents a small sample
        on ideal terrain.
  \item Supplementary video footage of all five locomotion trials
        is available at
        \url{https://github.com/romil0001/Sim-to-Real_balance_control_of_an_8-DoF_quadruped_using_ROS-2_and_embedded_complementary_filtering}.
\end{itemize}
Integrating the balance controller with the gait planner using
a switching or admittance-control interface is the primary
objective of continuing work (Section~\ref{sec:conc}).
\begin{figure}[!b]
  \centering
  \includegraphics[width=\columnwidth]{fig15.png}
  \caption{Open-loop trot gait trial on level ground.
    \textit{Top}: Bézier foot trajectories for diagonal pairs
    RF+LR (blue) and LF+RR (orange); step length $L_s=80\,\text{mm}$,
    foot clearance $H_\text{lift}=40\,\text{mm}$.
    \textit{Bottom}: Representative IMU pitch trace over a 2\,m
    forward-locomotion trial at $0.18\,\text{m\,s}^{-1}$; shaded
    bands indicate swing phases (PID disabled); stance-phase
    deviations remain ${<}1.5\dg$ with PID active.
    Supplementary video: \url{https://github.com/romil0001/Sim-to-Real_balance_control_of_an_8-DoF_quadruped_using_ROS-2_and_embedded_complementary_filtering}.}
  \label{fig:gait}
\end{figure}
% ============================================================
\section{Conclusion and Future Work}
\label{sec:conc}
% [M2][M7] Validation scoped to pitch/roll; locomotion framing qualified.
This paper presented the Rombotics Q8 as a \emph{platform
contribution}: a fully 3D-printed, open-source quadruped for
under \pounds250 ($\approx$\$315) that provides a reproducible,
independently pitch/roll validated baseline for low-cost
legged-robot research.
The platform achieves: $<1\dg$ static balance error;
$<0.5\,\text{s}$ recovery from $\pm10\dg$ disturbances;
\SI{80}{\milli\second} pitch/roll step response;
near-linear drift ($<0.003\dg$/min) over 60 minutes
characterised by independent inclinometer for pitch/roll
($n=5$ runs; preliminary result);
\SI{10}{\milli\second} control loops with $\sigma=0.23\,\text{ms}$
jitter over \num{10000} iterations;
a static stability margin of $\pm49$--$53\dg$ characterised by
quasi-static ZMP analysis; and a directed-locomotion capability
baseline at $0.18\,\text{m/s}$ on level ground with 71\%
stance-deviation reduction from closed-loop balance
[$t(6){=}9.0$, $p{<}0.001$].
The per-axis PID tuning strategy, the geometrically derived
diagonal-differential yaw actuation law, the PWM-space
tibia-coupling approximation (analytically bounded foot error
$<\SI{7}{\milli\metre}$), and the comprehensive hardware
characterisation are intended as reproducible engineering
reference points for teams building at this cost tier.
The platform is designed to enable comparative benchmarking:
its open-source firmware, documented timing budget, and
independent sensor validation provide the infrastructure
for future RL policy deployment, sim-to-real transfer,
and adaptive-gain studies without requiring costly hardware upgrades
\cite{bellicoso2018advances}.
\noindent Future work will address:
\begin{enumerate}
\item \textbf{Drift validation}: increase long-duration drift
      trials to $n \geq 10$ to achieve full statistical power
      for the drift-rate specification.
\item \textbf{Multi-surface testing}: characterise balance and
      gait performance on carpet, inclined plywood ($5\dg$,
      $10\dg$), and outdoor paving to bound surface-dependent
      performance.
\item \textbf{Closed-loop gait}: integrating the balance
      controller with the trot state machine using a
      switching MPC or admittance-control interface.
\item \textbf{RL-based locomotion}: using Isaac Gym
      \cite{rudin2022learning} to train sim-to-real locomotion
      policies deployable on the ESP32 via TensorFlow Lite
      Micro.
\item \textbf{Adaptive gains}: scheduling PID gains as a
      function of estimated payload or surface inclination
      from barometric altitude change.
\item \textbf{Yaw heading}: adding a magnetometer (HMC5883L)
      or monocular visual odometry to close the yaw estimation
      and validation loop.
\item \textbf{Wi-Fi telemetry}: exploiting the ESP32's onboard
      Wi-Fi for real-time parameter tuning and state logging.
\item \textbf{Digital twin}: deploying the MATLAB/Simulink
      model of the Q8 as a Gazebo digital twin for faster
      control parameter iteration.
\item \textbf{Thermal characterisation}: systematic servo
      temperature profiling under sustained balance corrections
      for high-ambient-temperature deployment qualification.
\end{enumerate}
% ── Acknowledgements ──────────────────────────────────────────
\section*{Acknowledgements}
The author thanks Dr.~Shabnam Sadeghi Esfahlani for component
provision, Wixey instrument loan, and weekly supervision;
Dr.~Hossein Bisheh for SolidWorks modelling guidance;
and Reece Sithanen for fabricating all printed components
at Anglia Ruskin University FabLab.
% ── Appendix ─────────────────────────────────────────────────
\appendices
\section{Notation Table}
\label{app:notation}
\begin{center}
\begin{tabular}{cl}
\toprule
$\phi,\theta,\psi$   & Roll, pitch, yaw Euler angles \\
$\Delta\vect{\Theta}$& Orientation offset from reference \\
$L_1, L_2$          & Femur, tibia link lengths (\SI{150}{\milli\metre}) \\
$\theta_1, \theta_2$ & Femur, tibia joint angles \\
$H_F, H_T$          & Femur, tibia home PWM counts \\
$\delta$             & Femur PWM offset from home \\
$K_{p,i}$           & Proportional gains (pitch/roll: 1.5; yaw: 1.2) \\
$K_{i,i}$           & Integral gains (pitch/roll: 0.1; yaw: 0.05) \\
$K_{d,i}$           & Derivative gains (pitch/roll: 0.2; yaw: 0.35) \\
$\alpha_s$          & Output smoothing factor (0.5) \\
$\alpha$            & Complementary filter weight (0.98) \\
$K_u$               & ZN ultimate gain (pitch/roll: 2.50; yaw: 2.00) \\
$T_u$               & ZN ultimate period (pitch/roll: 0.42\,s; yaw: 0.58\,s) \\
$h_\text{CoM}$      & Centre-of-mass height above foot plane \\
$\text{SSM}$        & Static stability margin \\
$L_s$               & Gait step length (\SI{80}{\milli\metre}) \\
$H_\text{lift}$     & Gait foot clearance (\SI{40}{\milli\metre}) \\
\bottomrule
\end{tabular}
\end{center}
\section{Firmware File Summary}
\label{app:firmware}
The firmware is organised across four files totalling
$\sim$\SI{64}{\kilo\byte} compiled (ESP32 flash):
\begin{itemize}
\item \texttt{MajorQuad\_v0\_3.ino} --- main control loop
      (Algorithm~\ref{alg:loop}), gait state machine,
      serial logging.
\item \texttt{ServoLimits.h} --- joint calibration constants
      (Table~\ref{tab:servolimits}).
\item \texttt{pid\_config.h} --- per-axis PID gain constants,
      smoother factor $\alpha_s$, ZN relay parameters.
\item \texttt{Helpers.h} --- \texttt{computeTibiaPWM()} function
      (Eq.~\ref{eq:tibiapwm}), Bézier step generator
      (Eq.~\ref{eq:bezier}).
\end{itemize}
Available at: \url{https://github.com/romil0001/Sim-to-Real_balance_control_of_an_8-DoF_quadruped_using_ROS-2_and_embedded_complementary_filtering}
% ── References ───────────────────────────────────────────────
\begin{thebibliography}{99}
\bibitem{biswal2021development}
P.~Biswal and P.\,K. Mohanty,
``Development of quadruped walking robots: A review,''
\textit{Ain Shams Engineering Journal},
vol.~12, no.~3, pp.~2017--2031, 2021.
\bibitem{pratt2000legged}
G.\,A. Pratt,
``Legged robots at MIT: What's new since Raibert?,''
\textit{IEEE Robotics \& Automation Magazine},
vol.~7, no.~3, pp.~15--19, Sep. 2000.
\bibitem{vukobratovic2004zmp}
M.~Vukobratovi\'{c} and B.~Borovac,
``Zero-moment point --- thirty-five years of its life,''
\textit{International Journal of Humanoid Robotics},
vol.~1, no.~1, pp.~157--173, 2004.
\bibitem{chai2022survey}
H.~Chai \textit{et al.},
``A survey of the development of quadruped robots,''
\textit{Biomimetic Intelligence and Robotics},
vol.~2, no.~1, p.~100029, 2022.
\bibitem{caron2015stability}
S.~Caron, Q.-C. Pham, and Y.~Nakamura,
``Stability of surface contacts for humanoid robots,''
in \textit{Proc. IEEE ICRA}, Seattle, 2015.
\bibitem{sabatini2011kalman}
A.\,M. Sabatini,
``Kalman-filter-based orientation determination using
inertial/magnetic sensors,''
\textit{Sensors}, vol.~11, no.~10, pp.~9182--9206, 2011.
\bibitem{vitali2021robust}
R.\,V. Vitali, R.\,S. McGinnis, and N.\,C. Perkins,
``Robust error-state Kalman filter for estimating IMU orientation,''
\textit{IEEE Sensors Journal},
vol.~21, no.~3, pp.~3561--3569, 2021.
\bibitem{madgwick2010efficient}
S.\,O.\,H. Madgwick, A.\,J.\,L. Harrison, and R.~Vaidyanathan,
``Estimation of IMU and MARG orientation using a gradient
descent algorithm,''
in \textit{Proc. IEEE Int. Conf. Rehabil. Robot.}, Zurich, 2011.
\bibitem{oz2022complementary}
M.\,N. \"{O}z, S.~Budak, E.~Kurnaz, and A.~Durdu,
``Orientation determination in IMU sensor with complementary
filter,''
\textit{Turkish Journal of Forecasting},
vol.~6, no.~1, pp.~34--39, 2022.
\bibitem{rahman2023irondogmini}
M.\,H. Rahman \textit{et al.},
``A dynamic approach to low-cost design, development, and
computational simulation of a 12DoF quadruped robot,''
\textit{Robotics}, vol.~12, no.~1, p.~28, 2023.
\bibitem{tirumala2019gait}
S.~Tirumala \textit{et al.},
``Gait library synthesis for quadruped robots via augmented
random search,''
\textit{arXiv:1912.12907}, 2019.
\bibitem{spotmicro2020}
D.~Kottenstette,
``Spot Micro: An open source quadruped robot,''
\textit{GitHub}, 2020.
[Online]. Available: \url{https://github.com/mike4192/spotMicro}
\bibitem{jia2018stability}
Y.~Jia \textit{et al.},
``Stability criterion for dynamic gaits of quadruped robot,''
\textit{Applied Sciences}, vol.~8, no.~12, p.~2381, 2018.
\bibitem{jahn2019modular}
U.~Jahn, C.~Wolff, and P.~Schulz,
``Concepts of a modular system architecture for distributed
robotic systems,''
\textit{Computers (Basel)}, vol.~8, no.~1, p.~25, 2019.
\bibitem{hutter2016anymal}
M.~Hutter \textit{et al.},
``ANYmal --- a highly mobile and dynamic quadrupedal robot,''
in \textit{Proc. IEEE/RSJ IROS}, Daejeon, 2016, pp.~38--44.
\bibitem{bledt2018cheetah}
G.~Bledt \textit{et al.},
``MIT Cheetah~3: Design and control of a robust, dynamic
quadruped robot,''
in \textit{Proc. IEEE/RSJ IROS}, Madrid, 2018, pp.~2245--2252.
\bibitem{dicarlo2018mpc}
J.~Di Carlo, P.\,M. Wensing, B.~Katz, G.~Bledt, and S.~Kim,
``Dynamic locomotion in the MIT Cheetah~3 through convex
model-predictive control,''
in \textit{Proc. IEEE/RSJ IROS}, Madrid, 2018, pp.~1--9.
\bibitem{gehring2016anymal}
C.~Gehring \textit{et al.},
``Practice makes perfect: An optimisation-based approach to
adaptive whole body control for quadrupeds,''
in \textit{Proc. IEEE ICRA}, Stockholm, 2016, pp.~1719--1726.
\bibitem{kumar2021rma}
A.~Kumar, Z.~Fu, D.~Pathak, and J.~Malik,
``RMA: Rapid motor adaptation for legged robots,''
in \textit{Proc. Robotics: Science and Systems (RSS)}, 2021.
\bibitem{lee2020learning}
J.~Lee \textit{et al.},
``Learning quadrupedal locomotion over challenging terrain,''
\textit{Science Robotics},
vol.~5, no.~47, 2020.
\bibitem{rudin2022learning}
N.~Rudin, D.~Hoeller, P.~Reist, and M.~Hutter,
``Learning to walk in minutes using massively parallel deep
reinforcement learning,''
in \textit{Proc. Conference on Robot Learning (CoRL)}, 2022.
\bibitem{kau2019stanford}
N.~Kau, A.~Schultz, N.~Ferrante, and P.~Slade,
``Stanford Doggo: An open-source, quasi-direct-drive
quadrupedal robot,''
in \textit{Proc. IEEE ICRA}, Montreal, 2019, pp.~6309--6315.
\bibitem{chen2021embedded}
X.~Chen, Y.~Liu, and J.~Wang,
``Real-time embedded control system design for multi-rotor
aerial platforms using ESP32 dual-core architecture,''
\textit{IEEE Access}, vol.~9, pp.~88\,241--88\,252, 2021.
\bibitem{liu2024sensing}
Y.~Liu \textit{et al.},
``A review of sensing technologies for indoor autonomous
mobile robots,''
\textit{Sensors (Basel)}, vol.~24, no.~4, p.~1222, 2024.
\bibitem{hwangbo2019learning}
J.~Hwangbo \textit{et al.},
``Learning agile and dynamic motor skills for legged robots,''
\textit{Science Robotics}, vol.~4, no.~26, 2019.
\bibitem{bellicoso2018advances}
C.\,D. Bellicoso \textit{et al.},
``Advances in real-world applications for legged robots,''
\textit{Journal of Field Robotics},
vol.~35, no.~8, pp.~1311--1326, 2018.
\bibitem{unitree2023go2}
Unitree Robotics,
``Go2: Consumer Quadruped Robot,''
\textit{Product Datasheet}, 2023.
[Online]. Available: \url{https://www.unitree.com/go2}
% [m7] New reference for ZN tuning methodology.
\bibitem{astrom1984automatic}
K.\,J. \r{A}str\"{o}m and T.~H\"{a}gglund,
``Automatic tuning of simple regulators with specifications on
phase and amplitude margins,''
\textit{Automatica}, vol.~20, no.~5, pp.~645--651, 1984.
\end{thebibliography}
\begin{IEEEbiography}[{\includegraphics[width=1in,height=1.25in,
  clip,keepaspectratio]{romil.png}}]{Romil Chetankumar Modi}
  received the MEng degree in Mechatronics and Robotics
  Engineering with Distinction from Anglia Ruskin University,
  Chelmsford, UK, in May 2025.
  His research interests include legged locomotion, real-time
  embedded control, computer vision, and IoT systems.
  He is the founder of the Rombotics open-source quadruped
  project (\url{https://rombotics.com}) and maintainer of the
  MajorQuad firmware
  (\url{https://github.com/romil0001/Sim-to-Real_balance_control_of_an_8-DoF_quadruped_using_ROS-2_and_embedded_complementary_filtering}).
\end{IEEEbiography}
\end{document}
